\documentclass[report]{itvest}

% Document-specific metadata
\title{Portfolio 1: Software Architecture}
\reportsubtitle{Case: Short Anonymised Case Name}
\author{Student Name} % group: \author{Name 1 \and Name 2 \and Name 3}
\date{\today}

% Portfolios use sections only, numbered 1, 2, ...
\renewcommand{\thesection}{\arabic{section}}

\begin{document}

\maketitle

\begin{checkit}
  \textbf{Assessment criteria} (delete this box before submitting):
  \begin{itemize}
    \item Scenarios are measurable and traceable to the business goal.
    \item Diagrams are consistent with the scenarios and the ADRs.
    \item ADRs name real alternatives and the trade-off that decided them.
    \item Debt risks and the AI Act classification are specific to the case,
      not a generic list.
  \end{itemize}
\end{checkit}

\section{Problem and Context}

State the business goal, the data available, and why ML is needed rather
than fixed rules. Describe the system at the level of components and data
types, and anonymise names where needed.

\begin{table}[h]
  \centering
  \begin{tabularx}{\textwidth}{lX}
    \toprule
    Stakeholder & Interest in the system \\
    \midrule
    \ldots & \ldots \\
    \ldots & \ldots \\
    \bottomrule
  \end{tabularx}
  \caption{Stakeholders.}
  \label{tab:stakeholders}
\end{table}

\section{Quality-Attribute Scenarios}
\label{sec:scenarios}

Give 3--5 scenarios in six-part form~\cite{bass_software_2021}. Say which
qualities matter most, which conflict, and what a wrong prediction costs.

\begin{checkit}
  At least one scenario is a \emph{model-quality scenario}, e.g.\ recall on
  a time-based hold-out set, with the cost of a missed case stated.
\end{checkit}

\begin{table}[h]
  \centering
  \footnotesize
  \begin{tabularx}{\textwidth}{lXXXXXX}
    \toprule
    ID & Source & Stimulus & Environment & Artefact & Response & Measure \\
    \midrule
    QA-1 & \ldots & \ldots & \ldots & \ldots & \ldots & \ldots \\
    QA-2 & \ldots & \ldots & \ldots & \ldots & \ldots & \ldots \\
    QA-3 & \ldots & \ldots & \ldots & \ldots & \ldots & \ldots \\
    \bottomrule
  \end{tabularx}
  \caption{Quality-attribute scenarios.}
  \label{tab:scenarios}
\end{table}

\section{Architecture}

Show the system as C4 context and container diagrams, with the AI and
non-AI parts marked. Name the existing systems it must integrate with
(e.g.\ ERP, historian, PLC/SCADA, data platform) and the network zones
it crosses.

\begin{figure}[h]
  \centering
  % Replace the box with: \includegraphics[width=\textwidth]{c4-context}
  \fbox{\parbox[c][5cm][c]{0.9\textwidth}{\centering
    Replace with the C4 context diagram.}}
  \caption{System context.}
  \label{fig:c4-context}
\end{figure}

\begin{figure}[h]
  \centering
  % Replace the box with: \includegraphics[width=\textwidth]{c4-container}
  \fbox{\parbox[c][7cm][c]{0.9\textwidth}{\centering
    Replace with the C4 container diagram, AI and non-AI parts marked.}}
  \caption{Containers.}
  \label{fig:c4-container}
\end{figure}

\section{Architecture Decisions}

Record at least two decisions. Each ADR names the scenarios from
Table~\ref{tab:scenarios} that drove it.

\subsection{ADR-1: Title of the Decision}

\begin{description}
  \item[Status:] Proposed / accepted.
  \item[Context:] The problem and the scenarios (QA-\ldots) that drive it.
  \item[Alternatives:] At least two real alternatives.
  \item[Decision:] The chosen alternative.
  \item[Consequences:] The trade-off that decided it, and what gets harder.
\end{description}

\subsection{ADR-2: Title of the Decision}

\begin{description}
  \item[Status:] \ldots
  \item[Context:] \ldots
  \item[Alternatives:] \ldots
  \item[Decision:] \ldots
  \item[Consequences:] \ldots
\end{description}

\section{EU AI Act Classification}

State the risk category of the system under the AI
Act~\cite{eu_ai_act_2024}, justify it with reference to the relevant
article or annex, and say what follows from it for the design.

\section{Technical-Debt Risks}

Name the debt the design will accumulate~\cite{sculley_hidden_2015}, and
how the design limits it.

\begin{table}[h]
  \centering
  \begin{tabularx}{\textwidth}{XXX}
    \toprule
    Debt risk & Where in the system & How the design limits it \\
    \midrule
    \ldots & \ldots & \ldots \\
    \ldots & \ldots & \ldots \\
    \bottomrule
  \end{tabularx}
  \caption{Technical-debt risks.}
  \label{tab:debt}
\end{table}

\section{LLM Considerations}

Only if your case uses an LLM or another foundation
model~\cite{bass_engineering_2025}: how prompts and model versions are
managed, how outputs are evaluated, and which guardrails apply. Otherwise
delete this section.

\section{AI Usage}

Briefly describe your use of AI tools on this assignment:

\begin{itemize}
  \item \emph{Code generation:} \ldots
  \item \emph{Grammar and writing:} \ldots
  \item \emph{Clarification of technical aspects:} \ldots
\end{itemize}

\bibliographystyle{IEEEtran}
\bibliography{bib}

\end{document}
